\NSPage{085}%
In particular, the covariance identity before the curl correction is
\NSFormula{NS-F-P085-001}{display}{none}
\[
  \mathcal C(W_0+\mathcal L\Sigma)=\varepsilon \mathcal T_{0,*}+\Sigma+\mathcal C(\mathcal L\Sigma).
\]
Thus \NSi{NS-F-P085-002}{\mathcal L} is a right inverse of
\NSi{NS-F-P085-003}{D\mathcal C(W_0)}. The quadratic remainder
\NSi{NS-F-P085-004}{\mathcal C(\mathcal L\Sigma)} is retained in the stress.
The same squared partition used for the leading covariance also assembles this linear inverse
across the slow boxes. Write \NSi{NS-F-P085-005}{W_{0,\beta}} and
\NSi{NS-F-P085-006}{\mathcal L_\beta} for the local constructions before multiplication
by the slow cutoff. For a real two-vector stress \NSi{NS-F-P085-007}{\sigma(r,z,t)}, independent of the angular and
auxiliary variables, whose normalized representatives \NSi{NS-F-P085-008}{Q^{2A}\sigma} are in
\NSi{NS-F-P085-009}{\mathsf M_\alpha}, set
\NSFormula{NS-F-P085-010}{display}{7.35}
\begin{equation}
\begin{aligned}
  W_{0,\mathrm{phys}}^{\mathrm{as}}
    &=\sum_{\beta=(\ell,a)} Q^{-A}\eta_\beta W_{0,\beta},\\
  \mathcal L_{\mathrm{phys}}^{\mathrm{as}}\sigma
    &=\sum_{\beta=(\ell,a)} Q^{-A}\eta_\beta\mathcal L_\beta(Q^{2A}\sigma).
\end{aligned}
\tag{7.35}\label{eq:7.35}
\end{equation}
Each summand uses the scale \NSi{NS-F-P085-011}{Q=2^{-\ell}} of its label
\NSi{NS-F-P085-012}{\beta=(\ell,a)}. The sums are locally finite for
\NSi{NS-F-P085-013}{q>0}. For overlapping slow supports, disjointness of the assigned rectangles, the squared
partitions, and the physical factor \NSi{NS-F-P085-014}{Q^{-2A}} give
\NSi{NS-F-P085-015}{\mathcal B(W_{0,\mathrm{phys}}^{\mathrm{as}},
\mathcal L_{\mathrm{phys}}^{\mathrm{as}}\sigma)=\sigma}. In a chart of scale
\NSi{NS-F-P085-016}{Q}, set \NSi{NS-F-P085-017}{\Sigma=Q^{2A}\sigma} and define the normalized representatives by
\NSFormula{NS-F-P085-018}{display}{none}
\[
  W_0^{\mathrm{as}}=Q^A W_{0,\mathrm{phys}}^{\mathrm{as}},
  \qquad
  \mathcal L^{\mathrm{as}}\Sigma
  =Q^A\mathcal L_{\mathrm{phys}}^{\mathrm{as}}(Q^{-2A}\Sigma).
\]
To test wave-class membership, collect all coefficients with the same label and harmonic, as
in Definition~\ref{def:6.5}. Bounded overlap of the slow cutoffs and the common-torus bounds give
\NSi{NS-F-P085-019}{W_0^{\mathrm{as}}\in \mathsf W_{1/2}} and
\NSi{NS-F-P085-020}{\mathcal L^{\mathrm{as}}\Sigma\in \mathsf W_{\alpha-1/2}}, and
\NSFormula{NS-F-P085-021}{display}{7.36}
\begin{equation}
  \mathcal B(W_0^{\mathrm{as}},\mathcal L^{\mathrm{as}}\Sigma)=\Sigma.
\tag{7.36}\label{eq:7.36}
\end{equation}
These definitions are consistent across overlapping charts. Auxiliary independence of the
input allows the scalar coefficients to be taken outside the covariance integral, as required
for the identity in Proposition~\ref{prop:7.6}.

\subsection{Exact curls, retained tails, and the covariance remainder.}\label{sec:7.4}
For the harmonic velocity
\NSi{NS-F-P085-022}{t_m e^{ikm\Phi}}, orthogonality to the phase gradient,
\NSi{NS-F-P085-023}{n_\Phi\cdot t_m=0}, cancels the leading term in its
divergence. Derivatives of the amplitude \NSi{NS-F-P085-024}{t_m} still contribute. Taking the full curl of the vector
potential \NSi{NS-F-P085-025}{A_m=C_m e^{ikm\Phi}} in \eqref{eq:7.38} supplies the corrected amplitude
\NSi{NS-F-P085-026}{t_m+r_m} needed for exact
incompressibility: the amplitude from differentiating the exponential is \NSi{NS-F-P085-027}{t_m}, while derivatives
of the potential coefficient \NSi{NS-F-P085-028}{C_m} and the cylindrical frame give
\NSi{NS-F-P085-029}{r_m}. In this subsection we
will estimate \NSi{NS-F-P085-030}{r_m} in Lemma~\ref{lem:7.7}, then the temporal cutoff residual in
\eqref{eq:7.40}, and finally the
covariance of a stress increment formed from the actual velocities in Corollary~\ref{cor:7.8}.

For a normalized potential \NSi{NS-F-P085-031}{A=(A_r,A_\theta,A_z)}, we define
\NSFormula{NS-F-P085-032}{display}{7.37}
\begin{equation}
  \operatorname{curl}_* A
  =\bigl(R^{-1}\partial_\theta A_z-D_zA_\theta,
         D_zA_r-D_rA_z,
         (D_r+R^{-1})A_\theta-R^{-1}\partial_\theta A_r\bigr).
\tag{7.37}\label{eq:7.37}
\end{equation}
The normalized derivatives include the chain-rule terms for evaluation at \eqref{eq:6.3}, so
\NSi{NS-F-P085-033}{\operatorname{curl}_* A} is
the normalized physical curl of the evaluated potential.

Fix a label, a nonzero integer \NSi{NS-F-P085-034}{m}, and a harmonic coefficient
\NSi{NS-F-P085-035}{t_m\in \mathsf W_\alpha} independent of
\NSi{NS-F-P085-036}{\theta}, with \NSi{NS-F-P085-037}{n_\Phi\cdot t_m=0}. We require
\NSi{NS-F-P085-038}{t_m} to be supported in the prescribed slow and transverse
supports and compactly inside the pulse interval, with smooth zero extension across all
these boundaries. We choose the potential so that the amplitude obtained when its curl
\NSPage{086}%
differentiates the exponential is \NSi{NS-F-P086-001}{t_m}. The potential coefficient
\NSi{NS-F-P086-002}{C_m} and the oscillating potential
\NSi{NS-F-P086-003}{A_m} are
\NSFormula{NS-F-P086-004}{display}{7.38}
\begin{equation}
  C_m=\frac{i\,n_\Phi\times t_m}{km|n_\Phi|^2},
  \qquad A_m=C_m e^{ikm\Phi},
  \qquad \operatorname{curl}_*A_m=(t_m+r_m)e^{ikm\Phi}.
\tag{7.38}\label{eq:7.38}
\end{equation}
The last identity defines \NSi{NS-F-P086-005}{r_m}, the difference between the curl amplitude and the prescribed
transverse amplitude. Since \NSi{NS-F-P086-006}{C_m} is independent of
\NSi{NS-F-P086-007}{\theta}, the remainder \NSi{NS-F-P086-008}{r_m} is
\NSFormula{NS-F-P086-009}{display}{none}
\[
  r_m=\bigl(-D_z(C_m)_\theta,
             D_z(C_m)_r-D_r(C_m)_z,
             (D_r+R^{-1})(C_m)_\theta\bigr).
\]

\begin{lemma}\label{lem:7.7}
For every \NSi{NS-F-P086-010}{\alpha\in\R} and nonzero integer
\NSi{NS-F-P086-011}{m}, let \NSi{NS-F-P086-012}{t_m\in \mathsf W_\alpha} satisfy
\NSi{NS-F-P086-013}{n_\Phi\cdot t_m=0}. Assume it
is supported in the prescribed slow and transverse supports and compactly inside the pulse interval,
with smooth zero extension across all these boundaries. The construction \eqref{eq:7.38} gives
\NSi{NS-F-P086-014}{C_m\in \mathsf W_{\alpha+1/2}}
and \NSi{NS-F-P086-015}{r_m\in \mathsf W_{\alpha+1/2-\kappa_s}}. The velocity
\NSi{NS-F-P086-016}{(t_m+r_m)e^{ikm\Phi}=\operatorname{curl}_*A_m} is exactly divergence-free for the
normalized operators and after evaluation at the phase map. Potential and velocity have smooth zero
extensions across their support boundaries. For the full amplitude
\NSi{NS-F-P086-017}{a_m=t_m+r_m}, the scalar product
with the phase normal satisfies
\NSFormula{NS-F-P086-018}{display}{7.39}
\begin{equation}
  ikm\,n_\Phi\cdot a_m
  =-\bigl((D_r+R^{-1})(a_m)_r+D_z(a_m)_z\bigr),
  \qquad n_\Phi\cdot a_m\in \mathsf W_{\alpha+1/2-\kappa_s}.
\tag{7.39}\label{eq:7.39}
\end{equation}
Physical potentials and pressures are obtained from chart ones by multiplication by
\NSi{NS-F-P086-019}{Q^{1/2-A}} and \NSi{NS-F-P086-020}{Q^{-2A}}, respectively.
\end{lemma}

\begin{proof}
The amplitude from differentiating the exponential in \NSi{NS-F-P086-021}{\operatorname{curl}_*A_m} is
\NSFormula{NS-F-P086-022}{display}{none}
\[
  ikm\,n_\Phi\times C_m
  =-\frac{n_\Phi\times(n_\Phi\times t_m)}{|n_\Phi|^2}
  =t_m-\frac{n_\Phi(n_\Phi\cdot t_m)}{|n_\Phi|^2}
  =t_m.
\]
Every remaining term in \NSi{NS-F-P086-023}{\operatorname{curl}_*A_m} differentiates the potential coefficient
\NSi{NS-F-P086-024}{C_m} or the cylindrical
frame. The factor \NSi{NS-F-P086-025}{(km)^{-1}} increases the
\NSi{NS-F-P086-026}{\varepsilon} exponent by \NSi{NS-F-P086-027}{1/2}, and applying
\NSi{NS-F-P086-028}{D_r} decreases it by at
most \NSi{NS-F-P086-029}{\kappa_s}, by \eqref{eq:6.32}. The multiplier
\NSi{NS-F-P086-030}{n_\Phi/|n_\Phi|^2} has polynomial coefficient bounds by Lemma~\ref{lem:7.1}.
These estimates prove the two class bounds with all derivatives of fixed order. The normalized
divergence is
\NSi{NS-F-P086-031}{\operatorname{div}_*a=(D_r+R^{-1})a_r+R^{-1}\partial_\theta a_\theta+D_z a_z}.
Since \NSi{NS-F-P086-032}{D_r(R^{-1})=-R^{-2}}, we have
\NSi{NS-F-P086-033}{(D_r+R^{-1})(R^{-1}f)=R^{-1}D_r f}. Substituting \eqref{eq:7.37} gives
\NSFormula{NS-F-P086-034}{display}{none}
\begin{align*}
  \operatorname{div}_*\operatorname{curl}_*A
   ={}&R^{-1}D_r\partial_\theta A_z-(D_r+R^{-1})D_zA_\theta\\
     &+R^{-1}D_z\partial_\theta A_r-R^{-1}\partial_\theta D_rA_z\\
     &+D_z(D_r+R^{-1})A_\theta-R^{-1}D_z\partial_\theta A_r=0,
\end{align*}
where the three pairs cancel by commutation. The chain rule preserves this identity after
evaluation at the phase map. Expanding the divergence of
\NSi{NS-F-P086-035}{a_m e^{ikm\Phi}}, whose amplitude
coefficient is independent of \NSi{NS-F-P086-036}{\theta}, gives \eqref{eq:7.39}; division by
\NSi{NS-F-P086-037}{km} and one application of \NSi{NS-F-P086-038}{D_r} give
the exponent improvement \NSi{NS-F-P086-039}{1/2-\kappa_s}. The support and smooth extension statements follow
from the coefficient bounds and cutoffs.
\end{proof}

To assemble real velocities, use conjugate pairs. For example,
\NSi{NS-F-P086-040}{t\cos(k\Phi)}, with real \NSi{NS-F-P086-041}{t},
has coefficients \NSi{NS-F-P086-042}{t/2} at harmonics
\NSi{NS-F-P086-043}{m=1,-1}. The corresponding potential and pressure
coefficients obey the same conjugacy condition. In particular, the pressure associated with a
real cosine velocity mode retains the factor \NSi{NS-F-P086-044}{i/(km)} in \eqref{eq:7.13}.

\paragraph{The temporal cutoff and its residual.}
For the solution of Proposition~\ref{prop:7.2}, we set
\NSi{NS-F-P086-045}{\widehat t_m=\psi t_m}
and \NSi{NS-F-P086-046}{\widehat\pi_m=\psi\pi_m}, and apply Lemma~\ref{lem:7.7} to
\NSi{NS-F-P086-047}{\widehat t_m}. We write \NSi{NS-F-P086-048}{\widehat r_m} for the curl remainder of this cutoff
amplitude. The solution before multiplication by \NSi{NS-F-P086-049}{\psi} may reach both pulse endpoints. The
\NSPage{087}%
cutoff makes the potential and pressure zero on an open neighborhood of each endpoint.
On the labelled rectangle \NSi{NS-F-P087-001}{0\leq v\leq L_s}, the product rule in \eqref{eq:7.5} gives
\NSFormula{NS-F-P087-002}{display}{7.40}
\begin{equation}
  \widehat t_m'+\mathcal K\widehat t_m+m^2d\widehat t_m
  +ikm n_\Phi\widehat\pi_m+f_m
  =(1-\psi)f_m+\psi't_m.
\tag{7.40}\label{eq:7.40}
\end{equation}
By \eqref{eq:6.16}, \NSi{NS-F-P087-003}{\psi=1\text{ on }|v-L_s/2|\leq L_s/5}, where both terms on the right vanish. Wherever
either term remains, \NSi{NS-F-P087-004}{|v-L_s/2|\geq L_s/5}, so \eqref{eq:7.16} gives an additional factor
\NSi{NS-F-P087-005}{S_*^C e^{-cS_*}} in every
fixed coefficient derivative estimate. In the iteration, \NSi{NS-F-P087-006}{f_m e^{ikm\Phi}} is an actual supported harmonic
source with smooth zero extension; thus this local identity describes its global uncancelled
remainder as well.

At a fixed correction stage, converting any fixed derivative estimate to physical variables
introduces a fixed power of \NSi{NS-F-P087-007}{Q^{-1}} and a polynomial in
\NSi{NS-F-P087-008}{S_*}. As \NSi{NS-F-P087-009}{S_*=\ell^2},
\NSi{NS-F-P087-010}{Q=2^{-\ell}}, and \NSi{NS-F-P087-011}{q\asymp Q},
for each fixed \NSi{NS-F-P087-012}{M,C,N} we have
\NSFormula{NS-F-P087-013}{display}{none}
\[
  q^{-N}Q^{-M}S_*^C e^{-cS_*}
  \leq C_N\exp\bigl(-c\ell^2+(M+N)\ell\log 2+2C\log\ell\bigr)
  \longrightarrow 0.
\]
Thus the resulting physical derivatives are \NSi{NS-F-P087-014}{O(q^N)} for every fixed
\NSi{NS-F-P087-015}{N}. These terms remain
separate additive flat residuals throughout subsequent corrections. The same argument
treats homogeneous cutoff tails. The pressure coefficient still satisfies
\NSi{NS-F-P087-016}{\widehat\pi_m\in \mathsf W_{\alpha+1/2}}, while
the exact velocity amplitude is
\NSi{NS-F-P087-017}{\widehat t_m+\widehat r_m}, with
\NSi{NS-F-P087-018}{\widehat r_m\in \mathsf W_{\alpha+1/2-\kappa_s}}.

\paragraph{Covariance of the actual velocities.}
The linear covariance inverse was defined at the fixed
primary wave. During iteration, its increment is added to a divergence-free wave field that
already contains earlier corrections, and the new increment has its own curl remainder. The
exact expansion below identifies the prescribed stress increment and retains three further
contributions: interaction with previous corrections, interaction with the new curl remainder,
and the square of the new velocity.

\begin{corollary}\label{cor:7.8}
Fix \NSi{NS-F-P087-019}{\alpha,\rho\in\R}, and use the normalized representatives of \eqref{eq:7.35}. Suppose
\NSi{NS-F-P087-020}{U} is a real
divergence-free wave field, including its curl remainders, such that
\NSFormula{NS-F-P087-021}{display}{none}
\[
  U=W_0^{\mathrm{as}}+E,
  \qquad U\in \mathsf W_{1/2},
  \qquad E\in \mathsf W_\rho.
\]
Let \NSi{NS-F-P087-022}{\Sigma(R,Z,T)\in \mathsf M_\alpha} be a real two-vector independent of the angular and auxiliary variables. Apply
Lemma~\ref{lem:7.7} to the signed amplitudes
\NSi{NS-F-P087-023}{\mathcal L^{\mathrm{as}}\Sigma}, obtaining the divergence-free increment
\NSi{NS-F-P087-024}{V=\mathcal L^{\mathrm{as}}\Sigma+R},
where \NSi{NS-F-P087-025}{R\in \mathsf W_{\alpha-\kappa_s}} is its curl remainder. Then
\NSFormula{NS-F-P087-026}{display}{7.41}
\begin{equation}
  \mathcal C(U+V)-\mathcal C(U)
  =\Sigma+\mathcal B(E,\mathcal L^{\mathrm{as}}\Sigma)
   +\mathcal B(U,R)+\mathcal C(V).
\tag{7.41}\label{eq:7.41}
\end{equation}
The three remainder terms satisfy
\NSFormula{NS-F-P087-027}{display}{7.42}
\begin{equation}
\begin{aligned}
  \mathcal B(E,\mathcal L^{\mathrm{as}}\Sigma)&\in \mathsf M_{\rho+\alpha-1/2},\\
  \mathcal B(U,R)&\in \mathsf M_{\alpha+1/2-\kappa_s},\\
  \mathcal C(V)&\in \mathsf M_{2\alpha-1}.
\end{aligned}
\tag{7.42}\label{eq:7.42}
\end{equation}
Taking the radial divergence of each remainder term decreases the
\NSi{NS-F-P087-028}{\varepsilon} exponent by at most a further
\NSi{NS-F-P087-029}{\kappa_s}.
\end{corollary}

\begin{proof}
We expand the quadratic map and use
\NSi{NS-F-P087-030}{\mathcal B(W_0^{\mathrm{as}},\mathcal L^{\mathrm{as}}\Sigma)=\Sigma}:
\NSFormula{NS-F-P087-031}{display}{none}
\begin{align*}
  \mathcal C(U+V)-\mathcal C(U)
    &=\mathcal B(U,V)+\mathcal C(V)\\
    &=\mathcal B(W_0^{\mathrm{as}},\mathcal L^{\mathrm{as}}\Sigma)
      +\mathcal B(E,\mathcal L^{\mathrm{as}}\Sigma)+\mathcal B(U,R)+\mathcal C(V)\\
    &=\Sigma+\mathcal B(E,\mathcal L^{\mathrm{as}}\Sigma)
      +\mathcal B(U,R)+\mathcal C(V).
\end{align*}
The bounds for \NSi{NS-F-P087-032}{\mathcal B(E,\mathcal L^{\mathrm{as}}\Sigma)} and
\NSi{NS-F-P087-033}{\mathcal B(U,R)} are products of the stated classes. Since
\NSi{NS-F-P087-034}{\kappa_s<1/2},
\NSi{NS-F-P087-035}{V\in \mathsf W_{\alpha-1/2}}, so
\NSi{NS-F-P087-036}{\mathcal C(V)\in \mathsf M_{2\alpha-1}}. This proves \eqref{eq:7.41} and \eqref{eq:7.42} without any positivity condition
\NSPage{088}%
on an accumulated stress. The radial-divergence assertion follows from \eqref{eq:6.32}; the zeroth-
order cylindrical terms cause no further exponent loss.
\end{proof}

\section{Compactly supported mean corrections}\label{sec:8}
The wave construction in Section~\ref{sec:7} provides pulses whose averaged covariance realizes the
leading stress target \eqref{eq:7.26}, and an inverse for the principal equation of each nonzero angular
harmonic (Proposition~\ref{prop:7.2}). The angular mean of the residual still requires correction. It
includes quadratic wave products and may depend on the auxiliary torus even though it is
independent of the physical angle (Proposition~\ref{prop:8.1}). In this section we construct pressure,
tangential stresses, and divergence-free velocity increments for this mean, with support
inside the active annulus (Proposition~\ref{prop:8.3} and Corollary~\ref{cor:8.5}).

There are two different inversion problems. For an auxiliary-averaged source, a compactly
supported solution \NSi{NS-F-P088-001}{\varphi} of
\NSi{NS-F-P088-002}{(D_r+e/R)\varphi=f}, with
\NSi{NS-F-P088-003}{e\in\{0,1,2\}}, requires
\NSi{NS-F-P088-004}{\int R^e f\,dR=0}; see
Lemma~\ref{lem:8.2}. For a source depending on the auxiliary torus, inversion of the fast time
derivative requires zero auxiliary average at each slow point (Lemma~\ref{lem:8.6}). We first derive the
mean momentum equations (Proposition~\ref{prop:8.1}) and construct the radial inverse (Lemma~\ref{lem:8.2}),
including the remainder needed when its source depends on the torus. We then apply it to
pressure and divergence-free axial increments (Proposition~\ref{prop:8.3}), and to tangential stresses
(Corollary~\ref{cor:8.5}). The integrated equations \eqref{eq:8.16} identify three scalar defects
\NSi{NS-F-P088-005}{\mathcal P,\mathcal J_\theta,\mathcal J_z}, defined
in Equations~\eqref{eq:8.12} and \eqref{eq:8.15}. A final linear map chooses five velocity coefficients to cancel
their linear changes while preserving two velocity moments (Lemma~\ref{lem:8.7}); we compute the
nonlinear changes explicitly in Lemma~\ref{lem:8.8}. The correction cycle in Proposition~\ref{prop:9.6} uses these
constructions after each wave update.

Unspecified radial integrals run from zero to infinity. We use the averaging and projection
conventions of Equations~\eqref{eq:3.7} and \eqref{eq:3.9}. Recall that
\NSi{NS-F-P088-006}{\langle\,\cdot\,\rangle_\theta} takes the angular average componentwise
in the cylindrical basis, whereas \NSi{NS-F-P088-007}{\langle\,\cdot\,\rangle_Y} takes Haar average on the auxiliary torus of
Lemma~\ref{lem:6.2}, and \NSi{NS-F-P088-008}{f^\circ=f-\langle f\rangle_Y}. In particular, an angularly invariant field may still depend
on that torus. The coefficient classes \NSi{NS-F-P088-009}{\mathsf M_\alpha} and
\NSi{NS-F-P088-010}{\mathsf S_\alpha} are those of Definition~\ref{def:6.4}; their estimates
hold for every fixed derivative order. Unless physical variables are displayed, we use the
normalized chart variables and operators of \eqref{eq:6.6}.

\subsection{Conservative momentum equations and integral constraints.}\label{sec:8.1}
In this subsection we
derive the mean momentum equations in divergence form (Proposition~\ref{prop:8.1}) and state the
two velocity moments \eqref{eq:8.2} that the corrections must preserve. The radial integrals of these
equations determine which compactly supported corrections are possible (Proposition~\ref{prop:8.4}
and Corollary~\ref{cor:8.5}); divergence form makes the boundary contributions explicit. The calculation
retains the quadratic products of the actual divergence-free field
\NSi{NS-F-P088-011}{w}, including its curl
remainders, as in \eqref{eq:8.1}. In a common chart with the normalized operators of \eqref{eq:6.6}, the actual
velocity has the decomposition
\NSFormula{NS-F-P088-012}{display}{8.1}
\begin{equation}
  u=(b+\beta,V+v,G+\gamma)+w,
  \qquad \langle w\rangle_\theta=0,
  \qquad W_{ab}=\langle w_aw_b\rangle_\theta.
\tag{8.1}\label{eq:8.1}
\end{equation}
Here \NSi{NS-F-P088-013}{(b,V,G)} is the slow base,
\NSi{NS-F-P088-014}{(\beta,v,\gamma)} is its angularly invariant correction, and
\NSi{NS-F-P088-015}{w} is the
sum of the divergence-free fields obtained by taking curls of the wave potentials. The indices
in \NSi{NS-F-P088-016}{W_{ab}} range over
\NSi{NS-F-P088-017}{a,b\in\{r,\theta,z\}}; thus \NSi{NS-F-P088-018}{W} includes every curl remainder. The correction
components and \NSi{NS-F-P088-019}{W} extend smoothly by zero outside the active shell. The mean correction is
divergence-free, meaning \NSi{NS-F-P088-020}{(D_r+R^{-1})\beta+D_z\gamma=0}, and we impose the two integral constraints
\NSFormula{NS-F-P088-021}{display}{8.2}
\begin{equation}
  M_\theta:=\int_0^\infty R^2\langle v\rangle_Y\,dR=0,
  \qquad
  M_z:=\int_0^\infty R\langle\gamma\rangle_Y\,dR=0.
\tag{8.2}\label{eq:8.2}
\end{equation}
\NSPage{089}%
These identities hold at every \NSi{NS-F-P089-001}{(Z,T)}. They are the zero angular-momentum moment and
zero axial-flux integral of the auxiliary-averaged correction, up to normalization; the base
is treated separately. Compactly supported azimuthal potentials preserve the axial-flux
constraint, and the finite-dimensional moment correction preserves both.

Let \NSi{NS-F-P089-002}{p_m} be the perturbation mean pressure, and put
\NSi{NS-F-P089-003}{\Delta_0=D_r^2+R^{-1}D_r+D_z^2} and
\NSi{NS-F-P089-004}{\Sigma_a=Q^{2A}\mathcal T_{\mathrm{phys},a}} for
\NSi{NS-F-P089-005}{a=\theta,z}, where the full residual stress tensor is that in \eqref{eq:5.41}.

\begin{proposition}\label{prop:8.1}
Let \NSi{NS-F-P089-006}{u} have the divergence-free decomposition \eqref{eq:8.1}, with slow base
\NSi{NS-F-P089-007}{(b,V,G)},
angularly invariant correction \NSi{NS-F-P089-008}{(\beta,v,\gamma)}, and
\NSi{NS-F-P089-009}{\langle w\rangle_\theta=0}. Let \NSi{NS-F-P089-010}{p_m} be its angularly invariant pressure
correction and let \NSi{NS-F-P089-011}{\Sigma_a=Q^{2A}\mathcal T_{\mathrm{phys},a}} be the normalized base stress components from \eqref{eq:5.41}. Retain the
base's flat residual as a separate additive error: all its derivatives of every fixed order vanish faster than
every power of \NSi{NS-F-P089-012}{q}. Then the angular mean of the remaining normalized residual is
\NSi{NS-F-P089-013}{(D_rp_m-g_r,E_\theta,E_z)}, where
\NSFormula{NS-F-P089-014}{display}{8.3}
\begin{equation}
\begin{aligned}
E_\theta={}&\mathfrak t_*v+(D_r+2/R)(bv+\beta V+\beta v+W_{r\theta})\\
 &+D_z(Gv+V\gamma+\gamma v+W_{z\theta})
   -\varepsilon(\Delta_0-R^{-2})v-(D_r+2/R)\Sigma_\theta,\\
E_z={}&\mathfrak t_*\gamma+(D_r+1/R)(b\gamma+\beta G+\beta\gamma+W_{rz})\\
 &+D_z(2G\gamma+\gamma^2+W_{zz}+p_m)-\varepsilon\Delta_0\gamma-(D_r+1/R)\Sigma_z,\\
g_r={}&-\mathfrak t_*\beta-(D_r+1/R)(2b\beta+\beta^2+W_{rr})\\
 &-D_z(b\gamma+G\beta+\beta\gamma+W_{zr})
   +(2Vv+v^2+W_{\theta\theta})/R+\varepsilon(\Delta_0-R^{-2})\beta.
\end{aligned}
\tag{8.3}\label{eq:8.3}
\end{equation}
\end{proposition}

\begin{proof}
Incompressibility puts the transport terms in divergence form. For cylindrical components of any divergence-free field, the three transport expressions are
\NSFormula{NS-F-P089-015}{display}{none}
\begin{align*}
  (u\cdot\nabla u)_r
    &=(D_r+1/R)u_r^2+R^{-1}\partial_\theta(u_\theta u_r)+D_z(u_z u_r)-u_\theta^2/R,\\
  (u\cdot\nabla u)_\theta
    &=(D_r+2/R)(u_r u_\theta)+R^{-1}\partial_\theta u_\theta^2+D_z(u_z u_\theta),\\
  (u\cdot\nabla u)_z
    &=(D_r+1/R)(u_r u_z)+R^{-1}\partial_\theta(u_\theta u_z)+D_z u_z^2.
\end{align*}
The rotation of the cylindrical frame contributes the centrifugal term in the first row and the
additional \NSi{NS-F-P089-016}{u_ru_\theta/R} in the second. Angular differentiation integrates to zero; products with
exactly one factor of \NSi{NS-F-P089-017}{w} do also. Subtracting the pure base terms therefore leaves the fluxes in
\eqref{eq:8.3}. The angular derivative terms in the vector Laplacian have zero angular average, leaving
\NSi{NS-F-P089-018}{\Delta_0-R^{-2}} in the radial and angular components and
\NSi{NS-F-P089-019}{\Delta_0} in the axial component. The residual
stress tensor contributes the two displayed radial divergences by \eqref{eq:5.41}. Excluded pulse tails
have nonzero angular harmonics and hence zero angular mean. The calculation retains the
full covariance of the divergence-free wave field, including the curl remainders.
\end{proof}

\subsection{A compactly supported primitive for the cylindrical weights.}\label{sec:8.2}
In this subsection we
construct a radial primitive with support in the active shell (Lemma~\ref{lem:8.2}). A direct radial
integral can remain nonzero beyond the source support. We will subtract its full radial
integral with a fixed cutoff and retain the resulting error. Since the radial derivative after
phase evaluation also differentiates the auxiliary variable \eqref{eq:6.4}, both integrals must follow
the corresponding path on the torus. We include that path in the physical definition below.

We define the radial operations first in physical variables, holding
\NSi{NS-F-P089-020}{(z,t)} and the independent angular variable fixed. We choose a slow cutoff
\NSi{NS-F-P089-021}{\chi_m(X)} that is zero on an inner collar,
one on an outer collar, and changes only in a fixed interior portion of the active shell. For a
smooth source \NSi{NS-F-P089-022}{g(r,z,t,Y)}, periodic in
\NSi{NS-F-P089-023}{Y} and smoothly extended by zero outside that shell,
\NSPage{090}%
set
\NSFormula{NS-F-P090-001}{display}{8.4}
\begin{equation}
\begin{aligned}
  \mathcal I g(r,Y)
    &=\int_0^r g\bigl(r',Y+((r')^{d_r}-r^{d_r})v_r\bigr)\,dr',\\
  \mathcal J g(r,Y)
    &=\int_0^\infty g\bigl(r',Y+((r')^{d_r}-r^{d_r})v_r\bigr)\,dr',\\
  \mathcal I_c g&=\mathcal I g-\chi_m\mathcal J g.
\end{aligned}
\tag{8.4}\label{eq:8.4}
\end{equation}
The shift in \NSi{NS-F-P090-002}{Y} follows the radial phase map while the physical
\NSi{NS-F-P090-003}{(z,t)} stay fixed. Below the
shell, \NSi{NS-F-P090-004}{\mathcal I g=0}; above it,
\NSi{NS-F-P090-005}{\mathcal I g=\mathcal J g}. Thus
\NSi{NS-F-P090-006}{\mathcal I_c g} vanishes for
\NSi{NS-F-P090-007}{X\leq X_a} and \NSi{NS-F-P090-008}{X\geq X_b}, with smooth
zero extension across both radial edges. The subtraction introduces an error supported
where \NSi{NS-F-P090-009}{\chi_m} varies, which we retain in the radial equation.

For the normalized operators, we use a common torus
\NSi{NS-F-P090-010}{y=Y_{i_0}} and put
\NSi{NS-F-P090-011}{M=\Lambda_g^{i_0}Q^{d_r}/2}. At
fixed \NSi{NS-F-P090-012}{(Z,T)}, the same symbols denote
\NSFormula{NS-F-P090-013}{display}{8.5}
\begin{equation}
\begin{aligned}
  \mathcal I g(R,y)
    &=\int_0^R g\bigl(R',y+M((R')^{d_r}-R^{d_r})v_r\bigr)\,dR',\\
  \mathcal J g(R,y)
    &=\int_0^\infty g\bigl(R',y+M((R')^{d_r}-R^{d_r})v_r\bigr)\,dR',\\
  \mathcal I_c g&=\mathcal I g-\chi_m\mathcal J g.
\end{aligned}
\tag{8.5}\label{eq:8.5}
\end{equation}
Here \NSi{NS-F-P090-014}{\chi_m=\chi_m(QR^2/(2q))} is the chart representative of the fixed physical cutoff. For
\NSi{NS-F-P090-015}{e\in\{0,1,2\}}, define
\NSFormula{NS-F-P090-016}{display}{8.6}
\begin{equation}
  \mathcal D_e=D_r+e/R,
  \qquad \mathsf T_e f=R^{-e}\mathcal I_c(R^e f),
  \qquad \mathcal A_e f=R^{-e}(\partial_R\chi_m)\mathcal J(R^e f).
\tag{8.6}\label{eq:8.6}
\end{equation}
At corresponding physical and chart points, with
\NSi{NS-F-P090-017}{r=\sqrt Q\,R}, the normalization is
\NSFormula{NS-F-P090-018}{display}{none}
\[
  f_*=Q^a f_{\mathrm{phys}},
  \qquad
  Q^{a-1/2}r^{-e}\mathcal I_{c,\mathrm{phys}}(r^e f_{\mathrm{phys}})
   =R^{-e}\mathcal I_c(R^e f_*)=\mathsf T_e f_*.
\]
Here \NSi{NS-F-P090-019}{\mathcal I_{c,\mathrm{phys}}} is the physical operator defined above. We call
\NSi{NS-F-P090-020}{\mathcal A_e f} the cutoff remainder
introduced by making the primitive compactly supported.

The inverse must preserve the decay order of its source. We also need its cutoff remainder
to be flat when the weighted auxiliary-mean integral vanishes. The next statement provides
both estimates.

\begin{lemma}\label{lem:8.2}
Fix a correction stage \NSi{NS-F-P090-021}{j}, an exponent
\NSi{NS-F-P090-022}{\alpha}, a band, its common torus
\NSi{NS-F-P090-023}{y=Y_{i_0}}, and
\NSi{NS-F-P090-024}{e\in\{0,1,2\}}. Let
\NSi{NS-F-P090-025}{f\in \mathsf M_\alpha} be a smooth scalar coefficient. In particular, its smooth zero extension
satisfies, for every fixed multi-index \NSi{NS-F-P090-026}{I} of derivatives in
\NSi{NS-F-P090-027}{(R,Z,T,y)},
\NSFormula{NS-F-P090-028}{display}{none}
\[
  |\partial^I f|\leq C_{j,I}\varepsilon^\alpha S_*^{B_{j,I}}\zeta\delta^{-B_{j,I}}.
\]
Then \NSi{NS-F-P090-029}{\mathsf T_e f} is supported in the same active radial shell and in the projection of the input support to
\NSi{NS-F-P090-030}{(Z,T)}, and belongs to \NSi{NS-F-P090-031}{\mathsf M_\alpha} under the primitive normalization just defined. Exactly,
\NSFormula{NS-F-P090-032}{display}{8.7}
\begin{equation}
  \mathcal D_e\mathsf T_e f=f-\mathcal A_e f.
\tag{8.7}\label{eq:8.7}
\end{equation}
If \NSi{NS-F-P090-033}{\int R^e\langle f\rangle_Y\,dR=0} at every
\NSi{NS-F-P090-034}{(Z,T)}, then \NSi{NS-F-P090-035}{\langle\mathcal A_e f\rangle_Y=0}, and for every fixed
\NSi{NS-F-P090-036}{I} and integer \NSi{NS-F-P090-037}{p\geq1},
\NSFormula{NS-F-P090-038}{display}{8.8}
\begin{equation}
  |\partial^I\mathcal A_e f|
  \leq C_{j,I,p}\varepsilon^{\alpha+p\kappa_s}S_*^{B_{j,I,p}}
         \zeta\delta^{-B_{j,I,p}}.
\tag{8.8}\label{eq:8.8}
\end{equation}
Only finitely many input derivatives are used for each
\NSi{NS-F-P090-039}{(I,p)}. Constants and powers may depend on
these indices, the stage, and the fixed profile, but not on the band, label, or point. If
\NSi{NS-F-P090-040}{f} is independent
of \NSi{NS-F-P090-041}{Y} and \NSi{NS-F-P090-042}{\int R^e f\,dR=0} at every slow point, the cutoff remainder is identically zero.
